% Generated draft: requires mathematical and source review.
\documentclass[11pt]{article}
\usepackage[margin=25mm]{geometry}
\usepackage{fontspec}
\setmainfont{Latin Modern Roman}
\usepackage{amsmath,amssymb,mathtools,unicode-math}
\setmathfont{Latin Modern Math}
\usepackage{enumitem,needspace,xurl,hyperref}
\setlength{\parindent}{0pt}
\setlength{\parskip}{5pt}
\setlength{\emergencystretch}{3em}
\newcommand{\N}{\mathbb{N}}
\newcommand{\Z}{\mathbb{Z}}
\newcommand{\Q}{\mathbb{Q}}
\newcommand{\R}{\mathbb{R}}
\newcommand{\C}{\mathbb{C}}
\newcommand{\F}{\mathbb{F}}
\newcommand{\A}{\mathbb{A}}
\newcommand{\PP}{\mathbb{P}}
\newcommand{\E}{\mathbb{E}}
\newcommand{\eps}{\varepsilon}
\newcommand{\dd}{\,\mathrm{d}}
\begin{document}
\section*{UnsolvedMath ID: 800292}
\textbf{Source problem code:} catalogue entry
\subsection{Definitions and mathematical statement}
Let \(G\) be a group. A left order on \(G\) is a strict total order \(\prec\) satisfying
\[
 a\prec b\ \Longrightarrow\ ca\prec cb
 \qquad(a,b,c\in G).
\]
Write \(\operatorname{LO}(G)\) for the set of all left orders on \(G\), and call \(G\) left-orderable when this set is nonempty. Give \(\operatorname{LO}(G)\) the topology generated by the sets
\[
 U_{a,b}=\{\prec\in\operatorname{LO}(G):a\prec b\},
 \qquad a,b\in G;
\]
thus a basic neighborhood specifies finitely many order comparisons. The conjugacy action of \(G\) on \(\operatorname{LO}(G)\) is defined by
\[
 a\mathrel{(g\cdot\prec)}b
 \quad\Longleftrightarrow\quad
 g^{-1}ag\prec g^{-1}bg
 \qquad(a,b,g\in G).
\]
This action is minimal if every orbit is dense, meaning
\[
 \overline{\{g\cdot\prec:g\in G\}}=\operatorname{LO}(G)
 \qquad\text{for every }\prec\in\operatorname{LO}(G).
\]
The question attributed to Navas asks whether there exists a group \(G\), generated by a finite subset and left-orderable, whose conjugacy action on \(\operatorname{LO}(G)\) is minimal.

\emph{Unresolved qualification.} The supplied statement does not require \(G\) to be nontrivial or infinite, and imposes no cardinality condition on \(\operatorname{LO}(G)\). These restrictions have not been added. Whether an additional restriction was intended requires source clarification.
\subsection{Short English statement}
Does there exist a finitely generated left-orderable group such that conjugating any one of its left orders produces a dense orbit in the space of all its left orders? The supplied wording leaves possible restrictions excluding trivial cases unspecified.
\subsection{Sources}
\begin{itemize}
\item Local UnsolvedMath registry: permanent ID \(800292\), source problem code \texttt{catalogue entry}, title \emph{Group Actions on 1-Manifolds --- Question 6}; collection: AMR Open Problem Lists. The supplied record attributes the question to Navas, with proposed year 2017, and labels its status as open. This is a supplied status label, not an independently verified literature assessment.
\item Navas, \emph{Group Actions on 1-Manifolds: A List of Very Concrete Open Questions} (2017), Question 6 (\#6). Supplied source URL: \url{https://arxiv.org/abs/1712.06462}. The supplied provenance records access on 2026-10-16 and extraction by \texttt{arxiv\_latex}; these are source-record provenance claims, not a claim of independent access or verification for this draft.
\end{itemize}
\end{document}
